\section{Smooth detector classes and complete-policy approximation}\label{sec:robust}
A polynomial mark is a design choice, not a consequence of collision
geometry alone. The geometry fixes the placement and collision masses.
For any specified smooth positive normalized circle density $q_a(t,y)$,
$t=(R-l)/(u-l)$, the same measured-coordinate apparatus can instead use
\[
 F_{r,a}(y)=\int_0^y q_a((r-l)/(u-l),v)\frac{dv}{2\pi},
 \qquad Y=F_{|z|,a}^{-1}(U).
\]
The known firmware evaluates its physical input $|z|$; the controller is
not given $R$. Positivity makes the CDF strictly increasing. This produces
the hit density $(2TR/a_0)q_a(t,y)$ while leaving the placement and no-hit
laws unchanged. Smooth positive firmware fits the same ideal local pointer
interpretation as Section~\ref{sec:lab}. The following results concern this
specified class, not all physically conceivable interventions.

\begin{theorem}[Uniform model error, optimal values, and transferred policies]\label{thm:model-stability}
Consider two apparatus families on the same report space and with the same
command and policy classes. Suppose their conditional next-report laws,
for every radius and command, differ in total variation by at most
$0\le e_0\le1$. Then for every common parameter-independent adaptive policy
and any common prior, the joint radius and length-$N$ report laws differ
by at most
\begin{equation}
 e_N=1-(1-e_0)^N\le\min\{1,Ne_0\}. \label{eq:coupled-model-error}
\end{equation}
If the common realized payoff lies in $[m_N,M_N]$, put
$\Omega_N=M_N-m_N$. Uniformly over any common policy subclass,
\begin{equation}
 |V_N-\widetilde V_N|\le\Omega_N e_N. \label{eq:model-value-error}
\end{equation}
A policy optimal for the second apparatus has regret at most
$2\Omega_N e_N$ in the first; an $\varepsilon$-optimal policy has regret
at most $2\Omega_N e_N+\varepsilon$. Applying a common raw-report censoring
channel does not increase $e_0$. No uniform posterior bound conditional
on an arbitrarily rare history is needed or asserted.
\end{theorem}
\begin{proof}
Fix the radius. While the two observed histories agree, the common policy
chooses the same command. Couple its next reports by their common part,
so their conditional agreement probability is at least $1-e_0$.
For the dominated families used here, the common part is explicitly the
minimum of the two densities, and the residual densities define a measurable
coupling. Use any coupling after the first mismatch. Induction bounds the
probability of agreement of all $N$ reports below by $(1-e_0)^N$, and the
coupling inequality proves \eqref{eq:coupled-model-error}. A common prior
is coupled identically at the start, giving the joint-law assertion.
A policy's independent random seed, if any, is coupled identically too.

A common censoring map and comparator seed take equal raw reports to equal
censored reports, proving contraction. For each fixed policy the payoff
is the same bounded measurable function of radius, report history, and
its induced commands and terminal decision in both experiments. Its
expectation changes by at most its oscillation times total variation.
Taking suprema over the same policy set proves \eqref{eq:model-value-error}.
Add and subtract the transferred policy's value in the second model to
obtain the factor-two regret bound, with the stated optimization error.
\end{proof}

For Section~\ref{sec:control} one may take
$m_N=G_-e^{-N|\lambda|B}$ and $M_N=G_+e^{N|\lambda|B}$; tighter
realized-payoff bounds are preferable when available. The linear telescoping
bound and the sinusoidal diagnostic below were proposed in the current
referee report. The next theorem supplies a positive implementable surrogate
for an entire smooth mark class, with a degree and regret certificate.

\begin{theorem}[Positive physical surrogates with degree--regret rates]\label{thm:smooth-surrogate}
Let the mode set be compact and let $q_a(t,y)$ be jointly continuous,
uniformly bounded above and below by positive constants, and normalized
by $\int q_a(t,y)dy/(2\pi)=1$. Assume smoothness of the firmware when a
smooth local pointer realization is invoked. For every integer $m\ge1$ define
\begin{equation}
 q_{a,m}(t,y)=\sum_{i=0}^m q_a(i/m,y)B_i^m(t). \label{eq:physical-Bernstein}
\end{equation}
This is a positive normalized mark family, implemented by the same
measured-coordinate inverse-CDF construction. Its complete raw kernel has
positive Bernstein coefficients at degree $d=m+1$; it keeps both physical
flag masses exactly. Every common-gate length-$N$ likelihood therefore
has an exact positive state with $dN+1$ coefficients, including failures.

Write $\|\cdot\|_1$ for integration in $y$ against $dy/(2\pi)$ and put
$p_{\max}=2T_+u/a_0$. If
$\|q_a(t,\cdot)-q_a(s,\cdot)\|_1\le L_1|t-s|$ uniformly, the raw
one-step model error is at most
\begin{equation}
 e_0\le p_{\max}\frac{L_1}{4\sqrt m}. \label{eq:Lipschitz-surrogate}
\end{equation}
If instead $q_a$ is twice continuously differentiable as an $L^1$-valued
map with $\sup_{a,t}\|\partial_t^2q_a(t,\cdot)\|_1\le L_2$, it is at most
\begin{equation}
 e_0\le p_{\max}\frac{L_2}{16m}. \label{eq:C2-surrogate}
\end{equation}
The right sides may be truncated at one. These bounds are uniform over
all gates, including endpoints, and all adaptive choices of mode.
An optimal surrogate policy has true-model regret at most
$2\Omega_N\{1-(1-e_0)^N\}$. In particular sufficient degrees for regret
at most $\delta>0$ are, respectively,
\begin{equation}
 m\ge\max\left\{1,\left(\frac{\Omega_N Np_{\max}L_1}{2\delta}\right)^2\right\},
 \qquad
 m\ge\max\left\{1,\frac{\Omega_N Np_{\max}L_2}{8\delta}\right\},
 \label{eq:degree-for-regret}
\end{equation}
with integer ceilings understood. These are sufficient implementation
bounds, not optimal approximation-dimension lower bounds.
\end{theorem}
\begin{proof}
The Bernstein basis is nonnegative and sums to one. Thus
\eqref{eq:physical-Bernstein} inherits the positive lower and upper bounds,
and integration in $y$ gives one exactly. Its samples are known firmware
values at prescribed radii, not measurements of the unknown radius.
Evaluating the polynomial at $(|z|-l)/(u-l)$ and applying its inverse CDF
implements it with the original prepared uniform seed.

The collision numerator is $2TR q_{a,m}$. Multiplication by the positive
linear polynomial $R=l+(u-l)t$ gives strictly positive degree-$(m+1)$
Bernstein coefficients by the product identity. The placement numerator
has positive degree-two coefficients $\pi(l^2,lu,u^2)$. The no-hit
numerator has degree-two coefficients bounded below by
$a_0-\pi u^2-2T_+u>0$. Elevation from degree two to $m+1\ge2$ preserves
positivity. Compactness gives uniform margins for every fixed $m$.
Their integrals normalize to $a_0$, so the whole raw kernel, its gate,
and its failure atom form an actual experiment. Normalizing every nonzero
failure factor as in Theorem~\ref{thm:boundary} includes gates with arbitrarily
rare failure. Positive convolution now gives the $dN+1$ coefficient state
and an attained finite-budget surrogate controller by the same compact
argument as in Section~\ref{sec:control} (with $d+1$ failure coefficients
in place of four).

For the approximation, let $S_m\sim\mathrm{Binomial}(m,t)$. Equation
\eqref{eq:physical-Bernstein} is $\E q_a(S_m/m,\cdot)$ in $L^1$.
The Lipschitz assumption and Cauchy--Schwarz give
\[
 \|q_{a,m}(t,\cdot)-q_a(t,\cdot)\|_1
 \le L_1\E|S_m/m-t|
 \le L_1\sqrt{t(1-t)/m}\le L_1/(2\sqrt m).
\]
For the second bound apply the Banach-valued Taylor formula with integral
remainder. Its linear term has zero expectation since $\E(S_m/m)=t$.
The remainder norm is bounded by $(L_2/2)|S_m/m-t|^2$, giving
\[
 \|q_{a,m}-q_a\|_1\le\frac{L_2 t(1-t)}{2m}\le\frac{L_2}{8m}.
\]
Total variation is half the $L^1$ norm. Only the hit component changes,
with probability $2TR/a_0\le p_{\max}$. This proves
\eqref{eq:Lipschitz-surrogate}--\eqref{eq:C2-surrogate}. Common censoring
contracts it. Theorem~\ref{thm:model-stability} supplies the policy and
regret statements. Its weaker $e_N\le Ne_0$ yields the sufficient degrees
\eqref{eq:degree-for-regret}.
\end{proof}

Given the one-step degree-$d$ coefficients and precomputed combinatorial
weights, multiplication at step $n$ uses $O(nd^2)$ positive arithmetic
operations; total filtering uses $O(N^2d^2)$ such operations and $dN+1$
stored coefficients. These counts exclude the computation of gate integrals,
firmware calibration, and global Bellman optimization. Combining the
surrogate regret with a separately certified optimization or lookup error
adds that error, rather than assuming it is zero in a numerical solver.
The payoff oscillation $\Omega_N$ in \eqref{eq:degree-for-regret} is explicit
and may grow exponentially with $N$ for a general multiplicative reward;
there is no hidden claim of a horizon-uniform $O(N/m)$ bound.

\subsection{A smooth physical perturbation that is not coefficient roundoff}
For $q_a=1+(\epsilon R^2+\delta\sin R)\cos(y-\theta)$, take $\delta$ small
enough to preserve strict positivity on $I$. The added hit numerator is
$2T\delta R\sin R\cos(y-\theta)$, whose fourth radius derivative is
\[
 2T\delta\{R\sin R-4\cos R\}\cos(y-\theta).
\]
It is not identically zero. The original exact degree-three cutoff does not
extend to this changed firmware. Nevertheless the exact raw one-step
variation distance is
\[
 \frac{2TR|\delta\sin R|}{a_0\pi},
\]
since the circle mean of $|\cos y|$ is $2/\pi$. This quantity controls
complete-policy value and regret by Theorem~\ref{thm:model-stability}.
Theorem~\ref{thm:smooth-surrogate} additionally constructs a positive
finite-degree approximation of the changed detector itself. None of these
value bounds asserts convergence of radius derivatives from total variation
alone; jet approximation would require its own derivative estimates.
The strong polynomial response theorem remains exact for each surrogate,
and for the original specified cubic detector.
